% Cover letter to the editor for the PeerJ submission.
% Standalone document (article class); builds independently of the manuscript.
\documentclass[11pt]{article}
\usepackage[margin=2.4cm]{geometry}
\usepackage[T1]{fontenc}
\usepackage{lmodern}
\usepackage{microtype}
\usepackage[colorlinks=true,linkcolor=black,urlcolor=blue]{hyperref}
\setlength{\parindent}{0pt}
\setlength{\parskip}{0.7em}
\pagestyle{empty}

\begin{document}

Pawel Kudzia\\
Motion, Mobility and Rehabilitation Lab\\
School of Exercise Science, Physical and Health Education\\
University of Victoria, Victoria, BC, Canada\\
pawel.kudzia5@gmail.com

\bigskip
Dear Editor,

We are pleased to submit our manuscript, ``Criterion validity of heel-mounted inertial sensors for stride-length estimation during single- and dual-task walking in older adults with and without cognitive impairment,'' for consideration as a Research Article in \textit{PeerJ Sports Medicine and Rehabilitation}.

Stride length is among the most informative gait markers of fall and cognitive-decline risk, yet it is also difficult to measure accurately outside the laboratory, particularly in the slow, variable gait of older adults. We validated a heel-mounted inertial measurement unit (IMU) pipeline against a GaitRite electronic walkway in 32 community-dwelling older adults, including 13 in the MoCA-defined cognitive impairment group, walking under single-task and three graded cognitive dual-task conditions. Across 10{,}472 matched strides from walks that passed synchronisation quality control, stride length agreed with the walkway to within 3.9~cm root-mean-square error (ICC 0.976), and 31 of 32 participants had a mean absolute error below 4~cm. Error remained low across cognitive load and in both cognitive-status groups.

This study extends heel-specific criterion validation to graded cognitive load in older adults across two MoCA-defined screening groups. Because the pipeline uses a clip-on heel sensor that needs no footwear modification and no per-participant calibration, it is well suited to further evaluation in clinical and free-living settings. We believe these findings will interest the \textit{PeerJ} readership in bioengineering, kinesiology, and geriatrics.

This manuscript is original, has not been published previously, and is not under consideration elsewhere. All authors have read and approved the submission and agree to its content. The authors declare no competing interests and received no specific funding for this work. The University of Victoria Human Research Ethics Board approved the study (protocol H22-00451), and all participants provided written informed consent. We have no preferred or excluded reviewers to suggest.

Thank you for considering our work. We look forward to your response.

\bigskip
Sincerely,

\medskip
Pawel Kudzia, on behalf of the co-authors (Markus von Hacht, Maryam Sheikhi, Drew Commandeur, Marc Klimstra, Sandra Hundza)\\
University of Victoria

\end{document}
